\cleardoublepage % or \clearpage for oneside

% ============================================
% APPENDICES (Optional)
% ============================================	
\appendix	
% Add commands directly to TOC file to custom prefic in ToC for appendix
\addtocontents{toc}{\protect\renewcommand{\protect\cftchappresnum}{Apx }}
\addtocontents{toc}{\protect\setlength{\protect\cftchapnumwidth}{1.5cm}}

\chapter{Additional Data}

Include supplementary material here. These are simple manual nomenclatures:

\section*{Symbols}
\begin{tabular}{ll}
$a$ & Acceleration \\
$v$ & Velocity \\
$m$ & Mass \\
$F$ & Force \\
\end{tabular}

\section*{Abbreviations}
\begin{tabular}{ll}
DNA & Deoxyribonucleic Acid \\
NASA & National Aeronautics and Space Administration \\
PhD & Doctor of Philosophy \\
\end{tabular}

\chapter{Code Listings}
\begin{spacing}{1.5}
    Include code or other technical details here. This is the main font (Times New Roman) - used for body text.
    \textsf{This is sans-serif (Arial) - used for headings.}	
    \texttt{This is monospace (Courier New) - used for code.}\par
    
    % Main text uses \setmainfont and \setmainhangulfont
    \noindent This is English in Times New Roman. 이것은 바탕체 한글입니다.\\
    % Sans-serif uses \setsansfont and \setsanshangulfont
    \textsf{This is English in Arial. 이것은 맑은 고딕 한글입니다.}\\
    % Monospace uses \setmonofont and \setmonohangulfont
    \texttt{This is code. 이것은 코딩 폰트입니다.}\par
    
  \begin{spacing}{2}
    \noindent{\tiny Tiny text (6pt)}\\
    {\scriptsize Scriptsize text (8pt)}\\
    {\footnotesize Footnotesize text (10pt)}\\
    {\small Small text (11pt)}\\
    \normalsize Normal text (12pt)\\
    {\large Large text (14pt)}\\
    {\Large Large text (17pt)}\\
    {\LARGE LARGE text (20pt)}\\
    {\huge Huge text (25pt)}\\
    {\Huge Huge text (25pt)}\par
  \end{spacing}
    
    % Custom sizes \fontsize{size}{baseline skip}\selectfont
    \noindent{\fontsize{15pt}{18pt}\selectfont Custom 15pt}\\
    {\fontsize{22pt}{26pt}\selectfont Custom 22pt}\\
    {\fontsize{30pt}{36pt}\selectfont Custom 30pt}\par
\end{spacing}

% ============================================
% ABSTRACT in HANGEUL
% ============================================	
\newpage	
\cleardoublepage
\phantomsection
\addcontentsline{toc}{chapter}{국문초록}
\vspace*{-1.5cm}
\begin{center}
{\LARGE\bfseries (\thesistitleKr)}
\\[3\baselineskip]
{\LARGE {\bfseries \thesisauthor }}
\\[2\baselineskip]
{\large 전남대학교 대학원 OO 학과\\[0.5\baselineskip] (지도교수 : \supervisorKr)}
\\[3\baselineskip]
\end{center}

{\Large \noindent {\bfseries (국문초록) }\par}
초록을 여기에 입력하세요. 논문의 요약은 일반적으로 150--350단어 분량으로, 문제, 방법론, 결과, 결론을 간결하게 설명해야 합니다.

\noindent\textbf{핵심어:} 첫 키워드, 두 번째 키워드, ...